\subsection{Independent exact source proof for the new Shift}
\label{sec:shift-independent}

An independent inspection of the resolved primitive program isolates the
changed Shift at source index 3660.  Write the nine logical registers as
$(C,T,W,A,D,H,X,Y,V)$ and suppose a command starts with
$T=t$, $W=p(a,b)$, $D=d$, $H=h$, and $A=X=Y=0$.
The candidate reader registers $C,V$ are arbitrary and remain unchanged.
The following statement concerns the resolved source program and uses no
Turing-machine arithmetic-box oracle.

The loop at 3660 clears $T$; each positive iteration consists of one
decrement and one jump, and its final zero test gives $2t+1$ instructions.
The loop at 3662 copies $H$ into both $T$ and $X$, consuming $H$;
its cost is $4h+1$.  The loop at 3666 transfers $X$ back into $H$,
at cost $3h+1$.  A final increment of $X$ and jump to index zero
produce
\[
 T=h,\qquad H=h,\qquad X=1,
\]
after exactly $2t+7h+5$ primitive instructions.

Indices zero through ten compare $H$ with $X$ and clear both on a
mismatch.  For $h=1$ they enter the halt instruction at index ten after
five further instructions.  For $h=0$ they reach index eleven after
five; for $h\ge2$ they reach eleven after $2h+5$.  Consequently
Shift enters halt exactly when $h=1$, after $2t+17$ source instructions.
This count stops on entry to the halt instruction; counting execution of
that instruction adds one.

For $h\ne1$, the routine at index eleven unpairs $W$ into
$X=a$ and $W=b$, with $Y=0$.  It then transfers $X$ into $H$ and
returns to the scheduler at index 25.  Here is a direct loop-count proof.
Put $k=a+b$ and $n=p(a,b)=k(k+1)/2+b$.  Each of the $n$ decrements
of $W$ costs five instructions apart from the carry transfer.
The successive carry transfers have lengths $1,2,\ldots,k$ and each
transferred unit costs three extra instructions. The final zero test,
the restoration of the residual $Y=b$ into $W$, and the transfer of
$X=a$ into $H$ cost $3a+3b+3$.  Thus the complete routine costs
\[
 5n+3\frac{k(k+1)}2+3k+3
   =4(a+b)^2+7a+12b+3.
\]
Its correctness follows from the corresponding triangular-counter
invariant: after a complete carry the pair of counters is $(X,Y)=(j,0)$,
and each subsequent noncarry step decrements $X$ and increments $Y$.
The endpoint $n=p(a,b)$ therefore leaves exactly $(a,b)$.

Combining these counts gives the all-natural statement
\[
 \begin{array}{ll}
 h=1:&\text{halt at index 10 after }2t+17\text{ instructions};\\[2pt]
 h\ne1:&\text{return to index 25 after }
 2t+9h+4(a+b)^2+7a+12b+13\text{ instructions},
 \end{array}
\]
with returning registers
\[
 (C,T,W,A,D,H,X,Y,V)=(C,h,b,0,d,a,0,0,V).
\]
All loops above decrease an explicitly displayed natural register; their
invariants and cost identities are quantifier-free polynomial equalities
and inequalities.  Primitive-recursive unpairing supplies $a,b$ for
arbitrary $W$, including $W=0$ with $a=b=0$.  This is a direct
quantifier-free-induction proof of the changed source command and its
runtime bound.  It does not independently establish the physical-table
simulation or the numerical-calculus calibration to PA.

The accompanying independent checker reads the exact primitive
instructions, pins each inspected block, and checks the formula on all
$0\le t,h,a,b<8$ (4096 cases).  Those finite cases audit the arithmetic
and indexing of the proof; they are not its induction argument.
